\documentclass[12pt,a4paper]{article}

\usepackage[utf8]{inputenc}
\usepackage[T1]{fontenc}
\usepackage[english]{babel}
\usepackage{amsmath, amssymb}
\usepackage{geometry}
\usepackage{enumitem}
\usepackage{fancyhdr}
\usepackage{xcolor}
\usepackage{graphicx}
\usepackage{hyperref}
\usepackage{cite}

% Page geometry configuration
\geometry{top=3cm, bottom=2.5cm, left=2.5cm, right=2.5cm}
\setlength{\headheight}{14.5pt}

% Header and Footer configuration
\pagestyle{fancy}
\fancyhf{}
\fancyhead[L]{\textbf{Olavo Barros (102014)}}
\fancyhead[R]{\textbf{INF 791 - Redes Complexas}}
\fancyfoot[C]{\thepage}
\renewcommand{\headrulewidth}{0.4pt}
\renewcommand{\footrulewidth}{0.4pt}

\begin{document}

% Custom Title Block
\begin{center}
    {\LARGE \textbf{LUT2Graph}}\\[0.5em]
    {\large \textbf{Intermediate Report --- Stage 1: Topic Definition}}\\[1em]
    \textcolor{gray}{\rule{\textwidth}{1pt}}\\[1em]
    \textbf{Student:} Olavo Barros \quad \textbf{ID:} 102014\\[0.5em]
    \textbf{Course:} INF 791 - Redes Complexas \\[0.5em]
    DPI - Universidade Federal de Viçosa (UFV)
\end{center}
\vspace{0.5em}


% -----------------------------------------------------------------------
\section*{(a) General Topic}

This project investigates \textbf{digital logic circuits} through the lens of
complex-network theory. Specifically, we focus on circuits that have been
synthesised by the Yosys open-source synthesis framework into networks of
\textbf{Look-Up Tables (LUTs)} --- the fundamental programmable logic elements
of FPGAs. Each LUT implements an arbitrary Boolean function of up to $k$
inputs by storing a truth table of $2^{k}$ binary entries.

The central idea is to map a synthesised netlist to a directed graph
$G = (V, E)$ and then apply complex-network analysis to characterise,
understand, and ultimately \textbf{prune} the circuit. Node and edge
annotations derived from information-theoretic measures (Boolean Bias and
Boolean Sensitivity) enrich the structural graph with functional meaning,
enabling pruning strategies that are grounded in both topology and logic.

% -----------------------------------------------------------------------
\section*{(b) Research Questions}

The project is guided by the following research questions:

\begin{enumerate}[label=\textbf{RQ\arabic*.}, leftmargin=*, itemsep=4pt]

    \item \textbf{Pruning via network metrics.}
    Can nodes with extreme Boolean Bias ($\beta \approx 0$ or
    $\beta \approx 1$) and edges with low Boolean Sensitivity
    ($\sigma_j \approx 0$) be safely removed without compromising the
    logical correctness of the circuit? What is the trade-off between
    graph reduction ratio and circuit fidelity?

    \item \textbf{Technique comparison.}
    How do classical complex-network pruning techniques (threshold-based
    edge removal, backbone extraction, community-based reduction) compare
    with LUT-specific pruning guided by Boolean Sensitivity?

\end{enumerate}

% -----------------------------------------------------------------------
\section*{(c) Dataset and Network Construction}

\paragraph{Dataset.}
The primary data source consists of synthesised netlists of benchmark circuits generated by TreeLUT \cite{khataei2025treelut} and Go-Fast \cite{barros2025go}--- a tool that optimizes XGBoost models to circuits for FPGA implementation. Was selected tree datasets, each containing 100 circuits of varying complexity (number of LUTs, depth, etc.), for the initial phase of the project. These circuits are representative of real-world applications and provide a rich testbed for evaluating pruning strategies.

\paragraph{Network construction.}
Each circuit is synthesised with \textbf{Yosys + ABC} using a 4-input LUT
technology mapping (\texttt{abc -lut 4; write\_json}). The resulting JSON
netlist is then parsed by the \texttt{LUTGraphBuilder} class (Python /
NetworkX) as follows:

\begin{itemize}[itemsep=2pt]
    \item \textbf{Nodes.} Each LUT cell becomes a node annotated with its
    Boolean Bias $\beta \in [0,1]$. Primary input and output ports are
    registered as boundary nodes.
    \item \textbf{Edges.} A directed edge $(u, v)$ is added for every
    connection from a driver signal $u$ to a driven LUT input $v$. The edge
    weight is the Boolean Sensitivity $\sigma_j \in [0,1]$ of that specific
    LUT input, quantifying how much the output changes when that input is
    toggled.
\end{itemize}

\bibliographystyle{plain}

\bibliography{references}

\end{document}